\documentclass[10pt,a4paper]{amsart}
\usepackage[margin=19mm]{geometry}
\usepackage{amsmath,amssymb,mathtools}
\usepackage[scr=rsfs,cal=euler]{mathalfa}
\usepackage{xfrac}
\usepackage[unicode,hidelinks,bookmarks=false]{hyperref}
\newcommand{\ie}{i.e.\ }
\newcommand{\cf}{cf.\ }
\newcommand{\resp}{respectively\ }
\newcommand{\Subsection}{\subsection}
\providecommand{\nobreakdash}{\nobreak\hbox{-}}
\setlength{\emergencystretch}{2em}
\multlinegap0pt
\usepackage[shortlabels]{enumitem}
\usepackage{amssymb}
\usepackage{bbm}
\usepackage{xcolor}
\newtheorem{cor}{Corollary}
\newtheorem{lem}{Lemma}
\usepackage{cleveref}
\crefname{cor}{Corollary}{Corollarys}
\crefname{lem}{Lemma}{Lemmas}
\newcommand{\N}{\mathbb{N}}
\renewcommand{\P}{\mathbb{P}}
\newcommand{\Q}{\mathbb{Q}}
\newcommand{\R}{\mathbb{R}}
\newcommand{\cC}{\mathcal{C}}
\newcommand{\cN}{\mathcal{N}}
\newcommand{\cS}{\mathcal{S}}
\renewcommand{\tt}[1]{{\mathtt{#1}}}
\title{The Airy line ensemble at the rough-smooth boundary}
\author{Sunil Chhita, Duncan Dauvergne, Thomas Finn}
\date{Source: arXiv:2602.24063v1 (CC BY 4.0)}
\begin{document}
\maketitle

\noindent\emph{Source and licence.} The Airy line ensemble at the rough-smooth boundary. Version arXiv:2602.24063v1 (CC BY 4.0), distributed by arXiv under the Creative Commons Attribution 4.0 International licence, http://creativecommons.org/licenses/by/4.0/, as recorded in the arXiv licence metadata for this exact version and shown on the abstract page https://arxiv.org/abs/2602.24063v1. Full licence notice: \texttt{licenses/arxiv-2602.24063-CC-BY-4.0.txt}.\par\smallskip

\noindent\textit{Editorial scope.} Counted unit: the article's lem L:QSE-pass (source0.tex lines 2147--2150). The article's complete original proof of it is delivered with it (source0.tex lines 2152--2237, one \texttt{proof} environment of 86 lines). No mathematics is written, repaired or re-derived: the statement and the proof are the article's own lines, in the article's own notation, and no retained line is substituted or re-flowed.  Retained uncounted as the source-local context the proof needs: lines 985--994; lines 1040--1051; lines 1242--1253; lines 1255--1266; lines 1288--1299; lines 1322--1333; lines 2031--2045; lines 2071--2092.  One declared adaptation: exactly 2 ancillary strings inside the retained spans are omitted (image-inclusion commands and, where the source repeats a label, the repeated label definitions) and no other line differs from the source; the figure environments, with their placement, captions and labels, are kept, and the package records the omitted strings in fidelity receipts bound to the affected bodies.  No retained line contains a citation, so the wrapper carries no bibliography and no bibliography entry.  The retained lines contain the article's own diagram code, which the wrapper typesets with the standard tikz packages; no image file is delivered and the package declares no asset dependency.  Editorial formatting only, in addition: an AMS \texttt{amsart} preamble that restates the article's own declarations and macros used by the retained lines, namely the environments \texttt{cor}, \texttt{lem} on separate counters, together with the standard packages those lines need.  The retained environments print in this document as Corollary 1, Lemma 1 in order of appearance, whereas the article prints the same items with the numbering of its own full document.  The complete native source of the article as distributed in the arXiv e-print for this exact version is bundled unchanged as \texttt{sources/195443/source0.tex}.
	\begin{cor}
		\label{C:pushforward}
		For $\tt{C} \in \{\tt{N}, \tt{S}\}$ and $a \in (0, 1)$, let $\Q_\tt{C, a, n}$ denote the pushforward measure of $\P_{a, n}$ under the map $F_\tt{C}$. Then $\Q_\tt{C, a, n}$ is a measure on dimer-compatible forests in $\tt{G(C)}$, where for a forest $F$ (oriented towards $\partial^\tt{x} \tt{C}$, $\Q_\tt{C, a, n}(F)$ is proportional to the weight of $F$, where edge weights on $\operatorname{Dir}_\tt{C}$ are given as follows:
		\begin{itemize}[nosep]
			\item $\tt{C} = \tt{N}$: NE/NW edges have weight $1$ and SE/SW edges weight $a^2$.
			\item $\tt{C} = \tt{S}$: SE/SW edges have weight $1$ and NE/NW edges weight $a^2$.


		\end{itemize}
	\end{cor}

	\begin{lem}
		\label{L:simple-sample}
		In this lemma, $\tt{G}(\tt{S})$ should be understood as a weighted directed graph, with south-biased weights given as in Corollary \ref{C:pushforward}, and $F \sim \Q_{\tt{S}, a, n}$.
		\begin{enumerate}[label=\arabic*.]
			\item Conditional on $\cS = \cS(F)$, $F \sim \tt{Forest}(\tt{G}(\tt{S}) \mid \cS)$.
			\item Let $v \in \tt{S}$, and let $\Pi^\cS(v) = (v = v_0, v_1, v_2, \dots, v_T)$ denote the path starting at $v$ in $F$, stopped at the first vertex $v_T \in \cS$. Then conditional on $\cS$, 
			$$
			(v_0, \dots, v_T) \sim \operatorname{LE}(X_v),
			$$
			where $X_v$ is a random walk on $\tt{G}(\tt{S})$ started at $v$, stopped when it hits $\cS$.
		\end{enumerate} 
	\end{lem}

	\begin{lem} (Global smooth coupling)
		\label{L:global-coupling-1} 
	There exist $a$-dependent constants $\gamma_a, d_a > 0$ such that the following holds. First define
		$$
		\Lambda_{\operatorname{sm}, 1} = \beta_n([-n^{5/6} , n^{5/6}] \times [n^{6/7}, n(|\xi_0(0)| + \gamma_a)). 
		$$
		Then
		$$
		d_{\operatorname{TV}}(\P_{a, n}|_{\Lambda_{\operatorname{sm}, 1}}, \P_a|_{\Lambda_{\operatorname{sm}, 1}}) \le 3 e^{-d_a n^{6/7}}.
		$$
		The same bound holds with $\Lambda_{\operatorname{sm}, 1}$ replaced by any of its rotations $R_{\pi/2} \Lambda_{\operatorname{sm}, 1}$, $R_{\pi} \Lambda_{\operatorname{sm}, 1}$, or $R_{3\pi/2} \Lambda_{\operatorname{sm}, 1}$, where $R_\theta$ denotes counterclockwise rotation by $\theta$.
	\end{lem}

	\begin{lem} (Mesoscopic smooth coupling)
		\label{L:global-coupling-2}
		Define the region
		$$
		\Lambda_{\operatorname{sm}, 2} = \beta_n([-n^{5/6} , n^{5/6}] \times [n^{1/3} \log^2 n, 2 n^{6/7}]). 
		$$
		Then there exists an $a$-dependent constant $d_a > 0$ such that
		$$
		d_{\operatorname{TV}}(\P_{a, n}|_{\Lambda_{\operatorname{sm}, 2}}, \P_a|_{\Lambda_{\operatorname{sm}, 2}}) \le 3 e^{- d_a \log^2 n}.
		$$
		The same bound holds with $\Lambda_{\operatorname{sm}, 2}$ replaced by any of its rotations $R_{\pi/2} \Lambda_{\operatorname{sm}, 2}$, $R_{\pi} \Lambda_{\operatorname{sm}, 2}$, or $R_{3\pi/2} \Lambda_{\operatorname{sm}, 2}$.
	\end{lem}

	\begin{lem}
		\label{L:lerw-estimate-basic}
		Let $X = (X_0, X_1, X_2, \dots)$ be a south-biased random walk on $\tt{\tilde S}$ started at $X_0 = (1, 0)$. For $\alpha \ge 2$, define the (roughly) parabolic region
		$$
		P^\alpha_\tt{S} = \{(x, y) \in \R^2 : y \le \alpha, |x| \le (\alpha + \log(y^- + 1))\sqrt{y^- + 1}\},
		$$
		where $y^- = \max(-y, 0)$.
		Then for an $a$-dependent constant $d_a > 0$, we have
		$$
		\P(\operatorname{Range}(X) \not\subset P^n_\tt{S}) \le 2 \exp (-d_a \alpha).
		$$
	\end{lem} 

	\begin{lem}
		\label{L:regular-paths}
		Let $F \sim \Q_{\tt{S}, a, n}$, let $\cS = \cS(F)$ be the backbone of $F$, and let $\Pi^\cS(v)$ denote the path in $F$ from a vertex $v$ stopped when it first hits the backbone, as in \cref{L:simple-sample}. Define the event
		$$
		\mathsf{Para}_\tt{S}(\alpha) = \{\operatorname{Range}(\Pi^\cS(v)) \subset P^\alpha_\tt{S} + v \text{ for all } v \in \tt{S}\}.
		$$
		Then
		$$
		\P(\mathsf{Para}_\tt{S}(\alpha) \mid \cS) \ge 1 - 2 n^2 \exp(-d_a \alpha).
		$$
        Here we are conditioning on the entire backbone of $F$, i.e.\ all south backbone paths.
	\end{lem}

	\begin{lem}
		\label{L:basic-path-ctrl}
		The following events hold with probability $1 - o(1)$ as $n \to \infty$.
		\begin{enumerate}[label=\arabic*.]
			\item The paths $\cS_i^\pm, \cN_i^\pm, i \ge n/4 + n^{1/4} \log^2 n$ do not enter the region 
			$$
			\tt{D^-} :=\tt{C}_{NE} \cup \tt{C}_{SW} \cup \tt{B}_{NE} \cup \tt{B}_{SW} \cup \tt{A}_\times.
			$$ 
			Similarly, the paths $\cS_i^\pm, \cN_i^\pm, i \le n/4 - n^{1/4} \log^2 n$ do not enter $\tt{D}^+ := \tt{C}_{NW} \cup \tt{C}_{SE} \cup \tt{B}_{NW} \cup \tt{B}_{SE} \cup \tt{A}_\times$. 
			\item The paths $\cS_i^\pm, \cN_i^\pm, |i - n/4| \le 2 n^{1/4} \log^2 n$ do not enter the set $\tt{C}_\times$.
			\item The paths $\cS_i^\pm, \cN_i^\pm, |i - n/4| \ge  n^{1/4} \log^4 n/2$ do not enter the set $\tt{A}_\times \cup \tt{B}_\times$.
			\item The split point $I_n$ satisfies $I_n - n/4 \in [-n^{1/4} \log^2 n-1, n^{1/4} \log^2 n]$. 
			
		\end{enumerate} 
	\end{lem}

	\begin{lem}
		\label{L:onion-control}
		For $r_1, r_2 \ge 1$ and $\ell \in \N$ let 
		$$
		\mathsf{Onion}(r_1, r_2, \ell)
		$$
		be the event where there exists a path $\cC_{i_0}^*:[0, t] \to \R$ for some $i_0 \in \{1, \dots, n/2\}$, $* \in \{+, -\}, \mathcal C \in \{\cN, \cS\}$ and path segments $\pi_1 = \cC_{i_0}^*|_{[s_1, t_1]}, \pi_2 = \cC_{i_0}^*|_{[s_2, t_2]}$ for some $t_1 \le s_2$ satisfying the following two conditions:
		\begin{itemize}
			\item There exists $\mathbf{v} \in \R^2$ such that $\pi_1, \pi_2 \subset B_{r_1, r_2}(\mathbf{v})$.
			\item The points $\pi_1(s_1), \pi_2(t_2)$ sit on the same boundary $\partial^\pm B_{r_1, r_2}(\mathbf{v})$ and the points $\pi_1(s_2), \pi_2(t_1)$ sit on the opposite boundary $\partial^\mp B_{r, R}(\mathbf{v})$.
			\item Let 
			$
			M
			\subset B_{r_1, r_2}(\mathbf{v})$ be the union of $\pi_1, \pi_2$ and the region in $B_{r_1, r_2}(\mathbf{v})$ bounded between these two paths. Then there are at most $\ell$ paths of the form $\cS_j^\pm$ or $\cN_j^\pm$ that enter $M$.
		\end{itemize}
		Informally, the event $\mathsf{Onion}(r_1, r_2, \ell)$ says that there exists an \textbf{onion} of width $r_1$, height $r_2$, and $\ell$ layers.
		Then assuming that $r_2 \ge 1$ and $r_1 \ge 10 \ell \sqrt{r_2} \log^2 n$, we have 
		\begin{equation}
			\label{E:Onion-r1-empty}
			\mathsf{Onion}(r_1, r_2, \ell) \cap \mathsf{Para}_\tt{S}(\log^2 n) \cap \mathsf{Para}_\tt{N}(\log^2 n) = \emptyset.
		\end{equation}
	\end{lem}

	\begin{lem}
		\label{L:QSE-pass}
		With probability $1 - o(1)$, the paths $\cS_i^-, i \le I$ do not enter $\tt{D}^+_0$. Similarly, the paths $\cS_i^-, i \ge I + 1, \cN_i^+, i \ge I + 1$ do not enter $\tt{D}^-_0$, and the paths $\cN_i^+, i \le I$ do not enter $\tt{D}^+_0$.
	\end{lem}

	\begin{proof}
		For the proof, we first define an event $E$ of probability $1 - o(1)$, and then show that the lemma holds deterministically on $E$. We will only verify the lemma for the paths $\cS_i^-, i \le I$, as the claim for the other three sets of paths follows by a symmetric argument.

      
		
		First, let $E_1$ be the probability $(1 - o(1))$-event where the four points in \cref{L:basic-path-ctrl} hold. Next, let $E_2$ be the event where for any path segment $\pi$ in the south forest $F_\tt{S}(D)$ started at a vertex $v$ and with $\pi \subset \tt{A}_\times$, we have
		\begin{equation}
			\label{E:pi-P}
			\pi \subset P_\tt{S}^{\log^2 n},
		\end{equation}
		and similarly for any path segment $\gamma$ in the north forest $F_\tt{N}(D)$ started at a vertex $v$ and with $\gamma \subset \tt{A}_\times$, we have	$\gamma \subset P_\tt{N}^{\log^2 n}$. The event $E_2$ has probability $1 - o(1)$ in the smooth phase by the standard estimate in \cref{L:lerw-estimate-basic} and a union bound. Therefore it also has probability $1 - o(1)$ in the Aztec diamond by the smooth phase couplings, \cref{L:global-coupling-1}, \ref{L:global-coupling-2}.
		We now define
		$$
		E = E_1 \cap E_2 \cap \mathsf{Para}_\tt{S}(\log^2 n) \cap \mathsf{Para}_\tt{N}(\log^2 n),
		$$
		which has probability $1 - o(1)$ by \cref{L:regular-paths},  \cref{L:basic-path-ctrl}, and the above discussion regarding $E_2$. In the remainder of the proof, we work on $E$. We will only prove the claim for $\tt{D}^+_0$, as the proof for $\tt{D}^-_0$ is symmetric. 	For the sake of contradiction, suppose that one of the paths $\cS_i^-, i \le I$ enters the set $\tt{D}^+_0$. By path ordering, the rightmost path $\cS_I^-$ must then also enter $\tt{D}^+_0$. 

          \begin{figure}
            \centering
            
            \caption{An illustration of the first part of the proof of \cref{L:QSE-pass}. For the the path $\tt{S}_I^-$ to cross over the line $L_2$, it must either cross through the sets $\tt{A}_\times$, or else use one of the corridors between $\tt{A}_\times$ and $\tt{C}_{SE}$ or $\tt{C}_{NW}$. The former is essentially rules out since south paths $\tt{A}_\times$  have a strong southward drift. The latter event is illustrated above, and results in an unfavourable $\tt{Onion}$ event.}
            \label{fig:lemma541}
        \end{figure}
        We first split $\tt{D}^+_0$ into two pieces,
		$$
		\tt{D}_1 = \{v \in \tt{D}^+_0 : v_1 + v_2 \ge - n^{5/6}/2\}, \qquad \tt{D}_2 = \tt{D}^+_0 \setminus \tt{D}_1.
		$$ 
		We first rule out that the paths $\cS_I^-$ enter $\tt{D}_1$. The argument is sketched in \cref{fig:lemma541}. Setting up some notation, define the lines
		$$
		L_i = \{(x, y) : x + y = - n^{5/6}/i \}, \quad i = 1, 2.
		$$
		Let $R$ denote the region bounded between the lines $L_1$ and $L_2$, and the sets $\tt{C}_{SE}, \tt{C}_{NW}$.
		
		For the sake of contradiction, suppose that one of the paths the paths $\cS_i^-, i \le I$ enters the set $\tt{E}_1$. By path ordering, the rightmost path $\cS_I^-$ must then also enter $\tt{E}_1$.	By Part (4) of \cref{L:basic-path-ctrl}, we have that $|I - n/4| \le n^{1/4} \log^2 n + 1$. Therefore by Part (2) of \cref{L:basic-path-ctrl}, the path $\cS_I^-$ cannot enter the sets $\tt{C}_{SE} \cup \tt{C}_{NW}$,	so the only way $\cS_I^-$ can enter $\tt{D}_1$ is if it crosses through the region $R$, first crossing the line $L_1$ and then later crossing the line $L_2$. Let $\pi_1 = \cS_I^-|_{[s, t]} \subset R$ denote the segment of $\cS_I^-$ given by taking the \textit{first} crossing of $R$ from $L_1$ to $L_2$. Let $v = \cS_I^-(s), w = \cS_I^-(t)$ denote the first and last points on $\pi_1$. We claim that
		\begin{equation}
			\label{E:n5610}
			w_1 - v_1 \ge n^{5/6}/100.
		\end{equation}
		Suppose that \eqref{E:n5610} fails. Then $w_2 - v_2 \ge 0.4 n^{5/6}$, since $L_1$ and $L_2$ are line segments parallel to the line $y = - x$, vertically shifted by $n^{5/6}/2$. On the other hand, the set $\tt{B}_{NW} \cap R$ is a thin strip contained in a set of the form
		$$
		\{(x, y) + v_{NW} : |x - y| \le n^{3/4}, |x+y| \le n^{5/6}/2\}
		$$
		for some $v_{NW} \in \R^2$. The same holds true of $\tt{B}_{SE} \cap R$ (with the vertex $v_{NW}$ replaced by a vertex $v_{SE}$). In particular, for $n$ large enough we have the estimate
		\begin{equation}
			\label{E:Cap-spread}
			\max \{|a_2 - a_2'| : a, a' \in \tt{B}_*\cap R \} < 0.39 n^{5/6}\}, \quad  * \in \{\text{NE}, \text{SW}\}.
		\end{equation}
		The negation of \eqref{E:n5610} together with \eqref{E:Cap-spread} implies that there exists a path $\pi' \subset \pi_1 \cap \tt{A}_\times$ starting and ending at vertices $a, a'$ and with $a'_2 - a_2 \ge n^{5/6}/100$. This contradicts the parabolic path estimate \eqref{E:pi-P}, and so \eqref{E:n5610} must hold. 

        \begin{figure}
            \centering
    
            \caption{The setup of the second part of the proof of \cref{L:QSE-pass}. Given that the path $\mathcal S_I^-$ does not enter $\tt{D}_1$, for it to enter $\tt{D}_2$ it must cross through the boundary strip $M$ in a way that contradicts a parabolic drift estimate. One example of this is given above.}
            \label{fig:lemma542}
        \end{figure}
		
		Now, since the path $\cS_I^-$ started on the south side of the Aztec diamond and finishes on the west side, there must exist a path $\pi_2 \subset \cS_I^-$ contained in the subset $S'$ of the strip $[v_1, w_1] \times \R$ which is bounded below by $\pi_1$:
		$$
		S' = \{(x, y) \in [v_1, w_1] \times \R : \text{ There exists } x \in \pi_1, v_2 \le y \},
		$$
		and which crosses $S'$ from right to left. Let $S''$ be the region in $S'$ bounded between $\pi_1$ and $\pi_2$. 
		
		By the ordering on paths, the only paths of the form $\cS^\pm_i$ that can enter the region $S''$ are the paths $\cS^-_i, i \le I$. Additionally, only those paths with index $i \ge n/4 - n^{1/4} \log^2 n$ can enter $S''$.	Indeed, the region $S''$ is contained in the union of the set $\tt{D}^+$ together with the component of the Aztec diamond that is connected to the north boundary. Therefore any path $\cS^-_i, i \le I$ that enters $S'$ must enter the region $\tt{D}^+$, and therefore must have index $i \ge n/4 - n^{1/4} \log^2 n$ by Part (1) of \cref{L:basic-path-ctrl}. 
		
		In summary, we have constructed a region in a strip $[v_1, w_1] \times \R$ of width at least $n^{5/6}/100$, bounded between two $\cS_i^-$ paths crossing from left to right and right to left, which contains at most $2 n^{1/4} \log^2 n \le 2 n^{1/4} \log^n$ $\cS_i$-paths. By $\cS$- and $\cN$-path interlacing, this region also contains at most $3n^{1/4} \log^2 n$ $\cN$-paths. That is, we are on the onion event
		$
		\mathsf{Onion}(n, n^{5/6}/100, 5n^{1/4}\log^2 n).
		$
		By \cref{L:onion-control}, this event cannot hold simultaneously with $E$. Having arrived at a contradiction, we can conclude that $\cS_I^-$ does not enter the set $\tt{D}_1$.
		
		To finish the proof, we show that $\cS_I^-$ does not enter $
        \tt{D}_2 = \tt{D}^+_0 \setminus \tt{D}_1$. The argument is illustrated in \cref{fig:lemma542}. Let $M$ be the closure of the set
		$$
		\{v \in \R^2 \setminus \tt{D}^+ : d(v, \tt{D}_2) \le n^{1/4}\}.
		$$
		Note that $M \subset \tt{A}_\times$, since when defining $\tt{D}_0^+$ we removed a boundary strip around the perimeter of $\tt{A}_\times$ in the region containing $\tt{D}_2$. Let $\partial^- M$ be the portion of the boundary of $M$ contained in $\tt{D}_2$, and let $\partial^+ M$ be the portion of the boundary not contained in $\tt{D}_0^-$.
		
		Since $\cS_I^-$ does not hit any vertex in $\tt{D}_1$, for $\cS_I^-$ to enter $\tt{D}_2$, there must be two disjoint path segments $\pi_1 = \cS_I^-|_{[s_1, t_1]}, \pi_2 = \cS_I^-|_{[s_2, t_2]}$ with $t_1 \le s_2$ contained in $M$, such that:
		\begin{itemize}
			\item $\cS_I^-(s_1), \cS_I^-(t_2) \in \partial^+ M$ and $\cS_I^-(s_2), \cS_I^-(t_1) \in \partial^- M$.
			\item The vertex $\cS_I^-(t_2)$ has a smaller polar coordinate in $(0, 2 \pi)$ than $\cS_I^-(s_1)$. We can guarantee this path ordering by choosing $[s_1, t_1]$ minimally, $[s_2, t_2]$ maximally, and using that $\cS_I^-$ moves from the south boundary to the west boundary of the Aztec diamond.
		\end{itemize}
		We will use these two points to derive a contradiction,  repeatedly using that $M \subset \tt{A}_\times$, so the estimate \eqref{E:pi-P} is in effect.
		
		First, by \eqref{E:pi-P} we must have $\cS_I^-(t_1) \in P_\tt{S}^{\log^2 n + 4} + \cS_I^-(s_1)$ (the $+ 4$ here is to accommodate the fact that the points $\cS_I^-(s_i), \cS_I^-(t_i)$ may not be vertices in $\tt{S}$). This forces the vertex $\cS_I^-(s_1)$ to lie on the portion of $\partial^{+} M$ where $y = x + n^{5/6}/2 + n^{1/4}$; call this part of the boundary $\partial^{++} M$. By the ordering in the second bullet, the point $\cS_I^-(t_2)$ must also lie on $\partial^{++} M$. However, again appealing to condition \eqref{E:pi-P}, we have that $\cS_I^-(t_2) \in P_\tt{S}^{\log^2 n + 4} + \cS_I^-(s_2)$. Since $\cS_I^-(s_2) \in \partial^- M$, this is a contradiction.
	\end{proof}

\end{document}
